\documentclass[10pt,a4paper]{amsart}
\usepackage[margin=19mm]{geometry}
\usepackage{amsmath,amssymb,mathtools}
\usepackage[scr=rsfs,cal=euler]{mathalfa}
\usepackage{xfrac}
\usepackage[unicode,hidelinks,bookmarks=false]{hyperref}
\newcommand{\ie}{i.e.\ }
\newcommand{\cf}{cf.\ }
\newcommand{\resp}{respectively\ }
\newcommand{\Subsection}{\subsection}
\providecommand{\nobreakdash}{\nobreak\hbox{-}}
\setlength{\emergencystretch}{2em}
\multlinegap0pt
\usepackage{amssymb}
\usepackage{enumerate}
\newtheorem{corollary}{Corollary}
\newtheorem{theorem}{Theorem}
\title{Rotational surfaces \\ with prescribed curvatures}
\author{Paula Carretero, Ildefonso Castro}
\date{Source: arXiv:2312.14672v2 (CC BY 4.0)}
\begin{document}
\maketitle

\noindent\emph{Source and licence.} Rotational surfaces \\ with prescribed curvatures. Version arXiv:2312.14672v2 (CC BY 4.0), distributed by arXiv under the Creative Commons Attribution 4.0 International licence, http://creativecommons.org/licenses/by/4.0/, as recorded in the arXiv licence metadata for this exact version and shown on the abstract page https://arxiv.org/abs/2312.14672v2. Full licence notice: \texttt{licenses/arxiv-2312.14672-CC-BY-4.0.txt}.\par\smallskip

\noindent\textit{Editorial scope.} Counted unit: the article's theorem Th:Kuhnel (source0.tex lines 264--266). The article's complete original proof of it is delivered with it (source0.tex lines 268--276, one \texttt{proof} environment of 9 lines). No mathematics is written, repaired or re-derived: the statement and the proof are the article's own lines, in the article's own notation, and no retained line is substituted or re-flowed.  Retained uncounted as the source-local context the proof needs: lines 200--203; lines 222--224; lines 258--260.  Nothing was omitted from any retained span, so the package carries no omission record.  The retained lines cite 1 works, whose entries are transcribed verbatim from the article's own bibliography source in the e-print and delivered with short editorial labels recorded in the item's reference mapping.  No retained line contains a figure environment, an image-inclusion command or drawing code, so no image is delivered and the package declares no asset dependency.  Editorial formatting only, in addition: an AMS \texttt{amsart} preamble that restates the article's own declarations and macros used by the retained lines, namely the environments \texttt{corollary}, \texttt{theorem} on separate counters, together with the standard packages those lines need.  The retained environments print in this document as Corollary 1, Theorem 1 in order of appearance, whereas the article prints the same items with the numbering of its own full document.  The complete native source of the article as distributed in the arXiv e-print for this exact version is bundled unchanged as \texttt{sources/151961/source0.tex}.
\begin{corollary}\cite[Corollary 2]{CC22}
	\label{cor:keysurfaces}
	Any rotational surface $S_\alpha$, with generatrix curve $\alpha = (x,z)$, is uniquely determined, up to $z$-translations, by the geometric linear momentum $\mathcal K=\mathcal K (x)$ of its generatrix curve, being $x$ non-constant. 
\end{corollary}

\begin{equation}\label{eq:prin curv}
	\kappa_1\equiv k_{\text m}= \mathcal K ' (x), \quad \kappa_2\equiv k_{\text p}= \frac{\mathcal K (x)}{x}.
\end{equation}

\begin{equation}\label{eq:K Kuhnel}
	\mathcal K (x)= \frac{x^q}{a^q}.
\end{equation}

	\begin{theorem}\label{Th:Kuhnel}
		The only rotational surfaces satisfying the linear Weingarten relation $k_{\text m} = q\,  k_{\text p}$, $q\neq 0$, are the plane and the Hopf-Kühnel surfaces $S_{\alpha_q}$.
	\end{theorem}

	\begin{proof}
		First, the plane trivially satisfies $k_{\text m} = q\,  k_{\text p}$, for any $q\neq 0$. Using \eqref{eq:K Kuhnel} in \eqref{eq:prin curv}, it is easy to check that $S_{\alpha_q}$ verifies precisely $k_{\text m} = q\,  k_{\text p}$.
		
		On the other hand, from \eqref{eq:prin curv} 
		the~linear Weingarten relation $k_{\text m}=q\,  k_{\text p}$ translates into the separable~o.d.e.
		$$\mathcal K ' (x)=q\, \mathcal K(x)/x. $$
		Its constant solution $\mathcal K \equiv 0$ leads to the plane.
		Its non-constant solution is given by $\mathcal K (x)= c \, x^q$, $c>0$. Taking $a=1/c^{1/q}$, we arrive at \eqref{eq:K Kuhnel}. The proof follows from Corollary \ref{cor:keysurfaces}.
	\end{proof}	

\begin{thebibliography}{99}
\bibitem[CC22]{CC22} P.~Carretero and I.~Castro. {\em A new approach to rotational Weingarten surfaces.} Mathematics {\bf 2022} 10(4), 578; https://doi.org/10.3390/math10040578
\end{thebibliography}
\end{document}
